\documentclass[12pt]{article}
\usepackage[utf8]{inputenc}
\usepackage[T1]{fontenc}
\usepackage[margin=1in]{geometry}
\usepackage{setspace}
\usepackage{xcolor}
\usepackage[colorlinks=true, urlcolor=blue!60!black, linkcolor=blue!60!black]{hyperref}
\setlength{\parindent}{0em}
\setlength{\parskip}{0.6em}
\renewcommand{\baselinestretch}{1.0}
\begin{document}
\begin{center}
{\LARGE \textit{Plessy v.\ Ferguson}, 163 U.S.\ 537 (1896)\par}
\vspace{0.5em}
{\large Case Brief --- Two-Page Edition\par}
\end{center}

\section{Memory Jogger}
``Separate but equal'' railway cars upheld as a reasonable police-power measure, over Justice Harlan's ``color-blind Constitution'' dissent (\href{https://supreme.justia.com/cases/federal/us/163/537/}{Justia}).

\section{Detailed Case Facts}
Louisiana's Separate Car Act of 1890 required ``equal but separate accommodations for the white, and colored races'' on intrastate railways. In a planned test case, Homer Plessy, one-eighth African in ancestry, sat in a whites-only coach in 1892, refused to move, and was arrested.

\section{Procedural History}
Judge John Howard Ferguson rejected Plessy's constitutional plea; the Louisiana Supreme Court denied relief. The U.S.\ Supreme Court affirmed.

\section{Judicial Votes}
Affirmed 7--1. Brown wrote for Fuller, Field, Gray, Shiras, White, and Peckham; Harlan dissented; Brewer did not participate.

\section{Holding}
State-mandated ``equal but separate'' accommodations on intrastate railways violate neither the Thirteenth Amendment nor the Fourteenth, which secures political but not ``social'' equality, if the law is a reasonable exercise of the police power grounded in local customs.

\section{Analysis of Opinions}
\textbf{Majority:} a law drawing ``merely a legal distinction'' between races is not servitude; segregated schools showed the practice's acceptance; and any badge of inferiority arose ``solely'' from how Black citizens chose to read the law. \textbf{Harlan:} ``Every one knows'' the Act's purpose was to exclude Black passengers from white coaches; it imposed a ``badge of servitude'' and would prove as pernicious as \textit{Dred Scott}.

\section{Examples: Future Applications}
\textbf{Same-Side:} segregated but nominally equal city buses. \textbf{Opposite-Side:} a whites-only law school with no equal alternative fails ``separate \textit{but equal}'' (\textit{Sweatt v.\ Painter}). \textbf{Fence-Sitter:} segregation on interstate trains, which collides with the Dormant Commerce Clause.

\section{Critique}
\textit{Plessy} belongs to the anti-canon; Charles Black exposed its willful blindness to segregation's purpose. Originalists debate whether the framers forbade segregation, and scholars dispute whether Harlan's dissent is color-blind or anti-caste.

\section{Key Quotations}
\begin{quote}
``But in view of the constitution, in the eye of the law, there is in this country no superior, dominant, ruling class of citizens. There is no caste here. Our constitution is color-blind, and neither knows nor tolerates classes among citizens. In respect of civil rights, all citizens are equal before the law.''
\end{quote}
\noindent --- Justice Harlan, dissenting, 163 U.S.\ at 559.

\end{document}
